\section{The Fisher--Rao marginal correction}
\label{sec:fr}

The Fisher--Rao (FR) correction is a birth--death / particle-reallocation step
that acts on the \emph{reaction-coordinate marginal} of the ensemble. It is
interleaved with the ABF dynamics~\eqref{eq:abf_sde} and is designed to drive the
empirical marginal $\phat_t(x)$ toward a target marginal $q_t(x)$, without
adding any drift to the $(x,y)$ dynamics.

\subsection{Empirical and target marginals}
At step $n$, after the Langevin proposal of all $N$ replicas, the empirical
reaction-coordinate marginal is estimated by a (reflected, edge-corrected)
kernel density / binned-smoothed estimate of bandwidth $\eta$,
\begin{equation}
  \phat_{n}(x) \;=\; \frac{1}{N}\sum_{i=1}^{N} K_\eta\!\bigl(x-X_n^i\bigr),
  \label{eq:phat}
\end{equation}
renormalised on the $x$-grid. The target $q_n(x)$ is one of three choices
(\cref{sec:targets}); all are nonnegative and normalised, $\int q_n=1$.

\subsection{Fisher--Rao score}
For each particle the FR \emph{score} is the pointwise log-ratio of the
empirical to the target marginal, with its ensemble mean removed:
\begin{equation}
  S_i \;=\;
  \log\frac{\phat_n(X^i_n)+\varepsilon}{q_n(X^i_n)+\varepsilon}
  \;-\;
  \underbrace{\int \phat_n(z)\,\log\frac{\phat_n(z)}{q_n(z)}\dd z}
  _{=\,\mathrm{KL}(\phat_n\,\Vert\,q_n)} ,
  \label{eq:score}
\end{equation}
where $\varepsilon>0$ is a small regulariser and the integral baseline equals
the Kullback--Leibler divergence $\mathrm{KL}(\phat_n\Vert q_n)$ (a grid
quadrature; it coincides with the Monte-Carlo baseline
$\tfrac1N\sum_k\log[\phat_n(X^k)/q_n(X^k)]$ up to discretisation). Subtracting
the baseline makes the score mean-zero, so the birth--death step below conserves
the expected particle number. The score is clipped to $[-S_{\max},S_{\max}]$
for stability. Its interpretation is direct:
\begin{itemize}
  \item $S_i>0$: $X^i$ lies where $\phat_n>q_n$, an \emph{over-represented}
        region (too many particles);
  \item $S_i<0$: $X^i$ lies where $\phat_n<q_n$, an \emph{under-represented}
        region (too few particles).
\end{itemize}

\subsection{Birth--death resampling}
The FR step approximates the Fisher--Rao gradient flow of
$\mathrm{KL}(\phat\Vert q)$ as a fixed-$N$ particle reallocation. Over a sub-step
of length $\Delta t$ and at effective rate $\gamma_n^{\mathrm{eff}}$
(\cref{sec:schedule}), each \emph{over-represented} particle ($S_i>0$) is marked
for death with probability
\begin{equation}
  \mathbb P(\text{$i$ dies}) \;=\; 1-\exp\!\bigl(-\gamma_n^{\mathrm{eff}}\,S_i\,\Delta t\bigr),
  \label{eq:death}
\end{equation}
and particles with $S_i<0$ supply clone sources. The WCA and metastability
backends implement this in the direct death-replacement form: killed slots are
filled by clones sampled from the under-represented pool, with weights increasing
with $|S_i|$. The entropic-bottleneck backend uses a vectorised but closely
related fixed-$N$ survivor/clone-pool variant. This distinction matters for exact
diagnostics, but not for the mechanism studied here: all variants move mass along
the reaction coordinate by copying whole configurations from under-represented
regions into over-represented ones.

Cloning copies the \emph{whole} configuration $(X^i,Y^i)$ --- the score depends
only on $x$, but the $y$-coordinate must travel with its $x$-coordinate, or
\cref{prop:cond_invariance} would be violated. The step is made deliberately
gentle by three safeguards: (i) a burn-in and soft ramp on the rate; (ii) score
clipping to $[-S_{\max},S_{\max}]$; and (iii) a hard cap $f_{\max}$ on the
fraction of particles replaced per event (\texttt{max\_event\_fraction}). The
numerical values of $S_{\max}$, $f_{\max}$ and the FR stride
\texttt{fr\_every} are case-specific and are listed in each case's design table.

\subsection{What marginal selection cannot do}
\label{sec:selection_invariance}
\Cref{prop:cond_invariance} says that an $x$-only \emph{bias} leaves the conditional law
$Y\mid X=x$ alone. The birth--death step is not a bias, so that proposition does not cover it,
and the natural worry is that resampling perturbs the conditional. The following says
something sharper, and in the opposite direction: in the mean-field limit the selection step
does not perturb the conditional at all, and therefore cannot repair it either.

\begin{proposition}[The selection flow leaves every conditional invariant]
\label{prop:sel_invariance}
Let $p_t(x,y)$ evolve under the mean-field birth--death flow generated by a score $S$ that
depends on the reaction coordinate alone,
\begin{equation}
  \partial_t p_t(x,y) \;=\; -\,\gamma\,\bigl[S(x)-\langle S\rangle_{p_t}\bigr]\,p_t(x,y),
  \qquad
  \langle S\rangle_{p_t}=\iint S(x)\,p_t(x,y)\dd x\dd y ,
  \label{eq:sel_flow}
\end{equation}
and write $p_t(x)=\int p_t(x,y)\dd y$ for the marginal and $p_t(y\mid x)=p_t(x,y)/p_t(x)$ for
the conditional, at any $x$ with $p_t(x)>0$. Then
\begin{equation}
  \partial_t\, p_t(x) \;=\; -\,\gamma\,\bigl[S(x)-\langle S\rangle_{p_t}\bigr]\,p_t(x),
  \qquad
  \partial_t\, p_t(y\mid x) \;=\; 0 .
  \label{eq:sel_conditional_frozen}
\end{equation}
\end{proposition}

\begin{proof}
Integrating~\eqref{eq:sel_flow} in $y$ and using that $S$ does not depend on $y$ gives the
first identity. For the second, differentiate the quotient and substitute both:
\[
  \partial_t \frac{p_t(x,y)}{p_t(x)}
  = \frac{\partial_t p_t(x,y)}{p_t(x)} - \frac{p_t(x,y)\,\partial_t p_t(x)}{p_t(x)^2}
  = \frac{-\gamma[S(x)-\langle S\rangle]\,p_t(x,y)}{p_t(x)}
    + \frac{\gamma[S(x)-\langle S\rangle]\,p_t(x,y)}{p_t(x)}
  = 0 . \qedhere
\]
\end{proof}

The factor $S(x)-\langle S\rangle$ multiplies $p_t(x,y)$ and $p_t(x)$ by the \emph{same}
number, so it cancels out of their ratio. Selection rescales the marginal and does nothing
else.

Three consequences organise the rest of the report.

\emph{First, the mechanism is bounded a priori.} Whatever a birth--death step buys, it is
bought by reweighting the marginal. It cannot be bought by improving the conditional average
that~\eqref{eq:cond_equal} says the mean-force estimate is built from. This is a stronger and
more useful statement than ``resampling is benign'': benign would permit a conditional
improvement, and \cref{prop:sel_invariance} forbids one.

\emph{Second, it is not a property of the Fisher--Rao score.} The proof uses only that $S$ is
a function of $x$. It therefore applies verbatim to the historical selection rules of
\cref{sec:wca_q1} --- the Chapter~6 rule $S\propto\partial_x^2 p/p$ and the count-balancing
rule $S\propto 1-p/\bar p$ --- which sit in exactly the same class. No score built on the
reaction coordinate alone can escape~\eqref{eq:sel_conditional_frozen}, so no such rule can
be distinguished from any other by its effect on conditionals. They can differ only in how
they redistribute the marginal, and \cref{sec:wca_q1} measures how much that distinction is
worth.

\emph{Third, it predicts a regime in which the method must fail.} If the binding difficulty
is that $p_t(y\mid x)$ has not equilibrated --- the transverse relaxation time
$\tau_\perp$ is long compared with the run --- then no marginal selection rule can help,
because none of them moves $p_t(y\mid x)$ at all. That is not a tuning failure to be fixed
with a better rate or a better target; it is structural. \Cref{sec:synthesis} makes it the
fourth regime of the map, and it is the reason the map now leads with a validity gate on the
ABF baseline rather than with a discovery test.

\begin{remark}[What the finite-$N$ ensemble does instead]
\label{rem:finite_n_conditional}
\Cref{prop:sel_invariance} is a mean-field statement, and the gap at finite $N$ runs the
unfavourable way. Cloning copies a whole configuration $(x,y)$, so a clone and its parent
share a $y$-coordinate: the conditional is preserved \emph{in expectation} while the number
of \emph{independent} $y$-samples behind each $x$-bin falls. The empirical conditional
therefore does not stay put --- it acquires resampling noise --- and the two safeguards that
matter are the ones that limit how fast lineages coalesce, namely the event-fraction cap
$f_{\max}$ and the requirement that clones decorrelate in time $\tau_\perp$ before the next
event. This is why the report tracks ancestor diversity alongside the free-energy error, and
why pushing the rate up eventually converts a marginal gain into a conditional loss.
\end{remark}

\subsection{Notation across cases}
\label{sec:fr_notation}
The same FR principle is applied to all three systems, but the implementation and
two parameter names differ slightly across backends. The birth--death \emph{rate}
is written $\gamma$ in the analytic two-dimensional cases
(\cref{sec:case_meta,sec:case_eb}) and \texttt{fr\_rate} in WCA
(\cref{sec:case_wca}); they play the same role in~\eqref{eq:death}. The
\emph{onset} of FR is parameterised as a burn-in \emph{fraction} $b$ in the
analytic cases (\eqref{eq:ramp}, $n_{\mathrm{burnin}}=\lfloor b\,
n_{\mathrm{steps}}\rfloor$) and as an absolute step count
\texttt{fr\_start\_steps} in WCA; the two are related by
$b\approx\texttt{fr\_start\_steps}/n_{\mathrm{steps}}$. Where a case uses a
collective variable $z(q)$ rather than $x$, read $x\mapsto z$ throughout this
section (\cref{rem:manybody}).

\subsection{Rate schedule}
\label{sec:schedule}
The FR rate is zero during an initial burn-in of $n_{\mathrm{burnin}}=
\lfloor b\,n_{\mathrm{steps}}\rfloor$ steps (burn-in fraction $b$), then ramps
softly to its maximum $\gamma$,
\begin{equation}
  \gamma_n^{\mathrm{eff}} \;=\;
  \gamma\,\Bigl(1-e^{-(n-n_{\mathrm{burnin}})/n_{\mathrm{ramp}}}\Bigr),
  \qquad n>n_{\mathrm{burnin}},
  \label{eq:ramp}
\end{equation}
and the birth--death step is applied every $\texttt{fr\_every}$ steps. The
burn-in lets ABF build a usable mean-force estimate (hence a usable target)
before any particles are moved; the ramp avoids a sudden shock to the ensemble.

\subsection{Target marginals: estimated, uniform, oracle}
\label{sec:targets}
The three targets differ only in which density $q_n$ the score~\eqref{eq:score}
is taken against. Recall $\Bt=\Fhat_n$ is the current ABF bias.

\paragraph{Estimated target (the practical method).}
A target free energy $\Fhat_n^{\mathrm{target}}$ is maintained independently of
the bias by an exponential moving average of the ABF estimate,
$\Fhat_{n+1}^{\mathrm{target}}=(1-r)\,\Fhat_n^{\mathrm{target}}+r\,\Fhat_{n+1}$,
and the target marginal is the \emph{biased} marginal that this target free
energy would induce under the current bias,
\begin{equation}
  q_n^{\mathrm{est}}(x) \;\propto\;
  \exp\!\Bigl[-\beta\bigl(\Fhat_n^{\mathrm{target}}(x)-\Bt(x)\bigr)\Bigr].
  \label{eq:q_est}
\end{equation}
This uses only quantities available at run time and is the method whose
practical value we assess.

\paragraph{Uniform target (a naive ablation).}
\begin{equation}
  q_n^{\mathrm{unif}}(x) \;=\; \mathrm{Unif}(x)\quad\text{on the $x$-grid.}
  \label{eq:q_unif}
\end{equation}
This tests whether simply \emph{flattening} the reaction-coordinate marginal is
enough. It is a target ablation, not the intended method; it ignores the fact
that, under a bias $\Bt\neq F$, the canonical biased marginal is not uniform.

\paragraph{Oracle target (a diagnostic control only).}
\begin{equation}
  q_n^{\mathrm{oracle}}(x) \;\propto\;
  \exp\!\Bigl[-\beta\bigl(\Fref(x)-\Bt(x)\bigr)\Bigr].
  \label{eq:q_oracle}
\end{equation}
\begin{remark}[The oracle target is not an algorithm]
\label{rem:oracle}
Equation~\eqref{eq:q_oracle} replaces the estimated target free energy by the
\emph{reference} free energy $\Fref$ --- exactly the unknown quantity the whole
procedure is trying to compute. The oracle target therefore \emph{cannot} be
evaluated in any real free-energy calculation. We include it solely as an
ideal-information diagnostic: it asks what the FR birth--death step would do if
the target marginal were anchored to the reference. This is not a universal upper
bound. If oracle FR improves on estimated FR in a given regime, the gap can be
read as target-shape headroom; if it ties or loses, the diagnostic instead says
that the limiting mechanism is not the online target estimate, and may even be a
mismatch between a fixed reference-shaped target and the transient ABF bias. The
practical comparison is always
\[
  \text{ABF only}\qquad\text{versus}\qquad\text{ABF}+\text{estimated-target FR}.
\]
\end{remark}

\subsection{The coupled ABF--FR step}
One outer iteration, given the state $\{(X^i_n,Y^i_n)\}$ and accumulators, is:
\begin{enumerate}
  \item update the ABF accumulators from the current particles and recompute
        $\Fphat_n,\Fhat_n=\Bt$ (\cref{sec:abf_estimator});
  \item propose the biased Langevin move~\eqref{eq:abf_sde} for all replicas;
  \item if FR is active at this step (past burn-in, on the \texttt{fr\_every}
        schedule): form $\phat_n$~\eqref{eq:phat} and the target $q_n$ ---
        one of~\eqref{eq:q_est},~\eqref{eq:q_unif},~\eqref{eq:q_oracle} ---
        compute the scores~\eqref{eq:score}, and apply the birth--death
        step~\eqref{eq:death};
  \item record diagnostics at the evaluation cadence.
\end{enumerate}
The ABF estimator sees the post-resampling ensemble at the next step, which is
the route by which improved $x$-coverage feeds back into a better mean-force
estimate. The remainder of the report quantifies whether, and how, this feedback
is beneficial.
